\documentclass[10pt,a4paper]{amsart}
\usepackage[margin=19mm]{geometry}
\usepackage{amsmath,amssymb,mathtools}
\usepackage[scr=rsfs,cal=euler]{mathalfa}
\usepackage{xfrac}
\usepackage[unicode,hidelinks,bookmarks=false]{hyperref}
\newcommand{\ie}{i.e.\ }
\newcommand{\cf}{cf.\ }
\newcommand{\resp}{respectively\ }
\newcommand{\Subsection}{\subsection}
\providecommand{\nobreakdash}{\nobreak\hbox{-}}
\setlength{\emergencystretch}{2em}
\multlinegap0pt
\usepackage{bm}
\usepackage{bbm}
\usepackage{dsfont}
\usepackage{xspace}
\usepackage{cancel}
\usepackage{mathtools}
\usepackage{amsthm}
\makeatletter
\@ifundefined{theorem}{\newtheorem{theorem}{Theorem}}{}
\@ifundefined{lemma}{\newtheorem{lemma}[theorem]{Lemma}}{}
\@ifundefined{proposition}{\newtheorem{proposition}[theorem]{Proposition}}{}
\makeatother
\title[On a Conjecture about Comparing the First and Second Zagreb ]{On a Conjecture about Comparing the First and Second Zagreb Indices of Graphs}
\author{Ali Ghalavand}
\date{Source: arXiv:2509.07434v1 (CC BY 4.0)}
\begin{document}
\maketitle

\noindent\emph{Source and licence.} On a Conjecture about Comparing the First and Second Zagreb Indices of Graphs. Version arXiv:2509.07434v1 (CC BY 4.0), distributed under the Creative Commons Attribution 4.0 International licence, as recorded in the arXiv licence metadata for this exact version and shown on the abstract page https://arxiv.org/abs/2509.07434v1. Full rights notice: licenses/arxiv-2509.07434-CC-BY-4.0.txt.\par\smallskip

\noindent\textit{Editorial scope.} Counted unit: the Proposition whose source label is \texttt{pro2} in the article (source0.tex lines 331--333, source label \texttt{pro2}), delivered together with the article's own complete original proof of it (source0.tex lines 334--366, one \texttt{proof} environment of 33 lines). No mathematics is written, repaired or re-derived: statement and proof are the article's own text line for line. Required context retained uncounted: lm1 (lines 174--177). Every definition, display and label of those lines is retained unchanged. Omitted from this package and not redistributed: 1--173, 178--330, 367--409. Nothing inside a retained excerpt is omitted or altered: every retained body is delivered exactly as the article's lines are written, so no omission receipt of any kind accompanies this package. Citations: the retained excerpts carry no citation command at all, so no bibliography entry of the article is transcribed and no reference is redistributed. Reference adaptation: none was needed and none was made; every reference inside the delivered unit resolves inside it, so no unresolved reference or citation is produced. Printed slips kept literal: no printed word, symbol, hypothesis or number of the counted statement or of its proof was altered. Layout and class: the article's own class is \texttt{article} with packages including amsmath, amsthm, amsfonts, amssymb, amscd, cite, tikz, color; the wrapper is the lane's pinned \texttt{amsart} editorial wrapper at 10pt on A4, which restates the theorem-like environments the retained text needs and adds the article's own macro definitions that the retained text needs (none), with extra wrapper packages (bm, bbm, dsfont, xspace, cancel, mathtools). The delivered numbering is the unit's own. Assets: no retained span contains a figure environment, a graphics-inclusion command, a PDF-page inclusion, a table or a TikZ picture; no image file is copied into this package, the package declares no asset dependency and no image file is redistributed with it. Licence: version arXiv:2509.07434v1 of this article is distributed under the Creative Commons Attribution 4.0 International licence, as recorded in the arXiv licence metadata for this exact version and shown on the abstract page https://arxiv.org/abs/2509.07434v1. A full rights notice is delivered at licenses/arxiv-2509.07434-CC-BY-4.0.txt. Characters: every retained excerpt is pure ASCII, so the delivered engine, XeLaTeX with its default Latin Modern text font, reports no missing character.
\begin{lemma}\label{lm1}
Let $ G_1 $ be a graph with a minimum degree of $ \delta(G_1) = 2 $ and a maximum degree of $ \Delta(G_1) = 5 $, such that $ m(G_1) = m_{2,5}(G_1) $. Additionally, let $ G_2 $ be a 3-regular graph. Consider a new graph $ G $ obtained by adding an edge between a vertex of degree $2$ in $ G_1 $ and a vertex of degree 3 in $ G_2 $. If $m(G_1)m(G_2)+360>81m(G_1)+80m(G_2)$, then 
 \[\frac{M_1(G)}{n(G)} > \frac{M_2(G)}{m(G)}.\]
\end{lemma} 

\begin{proposition}\label{pro2}
Let \( k \) be a positive integer, and let \( H \) be a 3-regular graph. Additionally, let \( G \) be the graph obtained from \( \Xi_{2k+1} \) and \( H \) by adding an edge between a vertex of degree 2 in \( \Xi_{2k+1} \) and an arbitrary vertex in \( H \). If the inequality $\frac{M_1(G)}{n(G)} > \frac{M_2(G)}{m(G)}$ holds true, then it follows that \( n(G) \geq 218 \).
\end{proposition}

\begin{proof}
By utilizing the structures of \( H \) and \( \Xi_{2k+1} \), along with the statements from the proof of Lemma \(\ref{lm1}\), we can observe the following relationships:
\begin{equation}\label{eqq1}
m(H)=\frac{3}{2}n(H),~m(\Xi_{2k+1})=\frac{10}{7}n(\Xi_{2k+1})=10(k+1).
\end{equation}
Furthermore, by applying the conclusions from the proof of Lemma \(\ref{lm1}\) again, we can see that if the inequality 
\[
\frac{M_1(G)}{n(G)} > \frac{M_2(G)}{m(G)}
\]
holds, it leads to the following inequality:
\begin{equation}\label{eqq2}
m(\Xi_{2k+1})m(H)+360-81m(\Xi_{2k+1})-80m(H)>0.
\end{equation}
By combining Equations \eqref{eqq1} and \eqref{eqq2}, we derive \( k \geq 8 \) and 
\[
15 n(H) + 360 - 810(k + 1) - 120 n(H) > 0.
\]
This simplifies to:
\[
n(H) > \frac{810k + 450}{15k - 105}.
\]
Since \( n(\Xi_{2k+1}) = 7(k + 1) \) and \( n(G) = n(H) + n(\Xi_{2k+1}) \), it follows that 
\[
n(G) > \frac{(7k + 19)(k - 1)}{k - 7}.
\]
Next, for \( 8 \leq k \leq 15 \), by utilizing Equations \(\eqref{eqq1}\) and \(\eqref{eqq2}\), and noting that \( n(H) \) is an even integer, we find that \( n(G) \geq 218 \). For \( k \geq 16 \), since the function \( f(k) = \frac{(7k + 19)(k - 1)}{k - 7} \) is increasing on the interval \( (16, \infty) \), we conclude that \( n(G) > f(16) > 218 \).
Additionally, if we assume that \( H \) is a 3-regular graph of order 106, and let \( G \) be the graph obtained from \( \Xi_{31} \) and \( H \) by adding an edge between a vertex of degree 2 in \( \Xi_{31} \) and an arbitrary vertex in \( H \), we can confirm that \( n(G) = 218 \) and that 
\[
\frac{M_1(G)}{n(G)} > \frac{M_2(G)}{m(G)},
\]
as desired.

\end{proof}

\end{document}
